\section{Introduction}\label{sec:intro}

In an Industrial Internet of Things (IIoT) deployment, thousands of sensors report readings that drive
physical processes: temperatures, pressures, valve states. A forged reading can open a valve or shut
down a production line, so every reading must be authenticated. Two engineering pressures shape how
this is done. First, the sensors are constrained and the network is shared, so the verifier should not
have to receive and check $n$ separate signatures; \emph{aggregate signatures}~\cite{boneh2003} let an
aggregator compress $n$ signatures into one object that the cloud checks with a single accept-or-reject
decision. Second, the key infrastructure should be light: a conventional PKI needs certificate storage,
distribution and revocation, while identity-based signatures~\cite{shamir1984} let the key authority
compute every user's private key (\emph{key escrow}). \emph{Certificate-based} cryptography~\cite{gentry2003}
sits between the two: each sensor generates its own key pair, a key generation centre (KGC) issues it a
certificate, and signing needs both, so neither the sensor's key nor the KGC alone suffices to sign.

That design goal only holds if the KGC cannot forge. The KGC holds a master secret key, and master keys
leak: the paper we study cites an incident in which an insider at South Africa's Postbank took a data
centre's master key and stole 3.2 million US dollars~\cite{qiao2023}. A certificate-based aggregate
signature (CBAS) scheme for IIoT must therefore remain unforgeable even against a \emph{malicious} KGC.

\paragraph{The paper under study.}
Qiao \emph{et al.}~\cite{qiao2023} (IEEE Systems Journal, 2023) examine the CBAS scheme of Verma
\emph{et al.}~\cite{verma2020} (CB-CAS, 2020). They show that a malicious KGC that has observed one valid
signature can forge signatures on arbitrary messages, and they propose a repaired scheme. They prove the
repair existentially unforgeable under the discrete-logarithm assumption in the random-oracle model, and
their cost tables present it as exactly as cheap as Verma's scheme.

\paragraph{Our questions.}
We set out to check the paper's claims independently and then to test the repaired scheme further:
\begin{enumerate}[label=\textbf{RQ\arabic*}]
\item Does the paper's attack work, and which change in the repair stops it?
\item Is the repaired scheme as cheap as the paper's tables state, and does a good implementation change
      the answer?
\item Is the paper's attack the weakest attack on Verma's scheme, and which of the repair's changes closes
      which attack?
\item What happens when the repaired scheme is deployed: with a faulty or malicious sensor in a batch,
      with a weak random-number generator, and against the extra security properties the paper claims?
\end{enumerate}

\paragraph{Contributions.}
We distinguish throughout between what the paper claims \prov{Paper}, what we re-ran \prov{Reproduced},
what we measured \prov{Experiment}, our design choices \prov{Implementation}, and results that are new to
this project \prov{Ours}. Known techniques that we only apply to this scheme are named as such.
\begin{enumerate}
\item An implementation of both schemes, faithful to the paper, over two independent prime-order
      groups, with domain-separated hashes, unambiguous encodings and instrumented operation counting, checked by 424
      tests that run on both groups (\cref{sec:method}). \prov{Implementation}
\item Reproduction of the paper's malicious-KGC forgery (8/8 trials) and of the repair's resistance to it
      (0/8), plus a controlled ablation showing that hashing the nonce commitment $T_i$ is the change that
      stops this attack (\cref{sec:repro}). \prov{Reproduced}, ablation \prov{Ours}
\item An audit of the paper's cost tables: the row given for the repaired scheme is identical to Verma's
      row, but its own verification equation costs $(2n+2)\Te+3n\Ta+3n\Th$, about twice the printed value
      ($\times1.99$ at $n=100$), confirmed by timing on two curves and two machines (\cref{sec:cost}).
      \prov{Ours}
\item A key-only universal forgery on Verma's scheme as described in the paper, which needs no observed
      signature and contradicts the paper's Table~II, and a $2\times2$ ablation showing that each of the
      repair's two hash bindings closes a different attack (\cref{sec:keyonly}). \prov{Ours}
\item A best-effort optimisation study showing that a better implementation does not close the cost gap:
      about $\times2.4$ before and after (\cref{sec:opt}). \prov{Experiment}
\item A deployment analysis: the cost of handling invalid signatures at the aggregator
      (\cref{sec:agg}), key recovery from a repeated signing nonce with a hedged-nonce fix
      (\cref{sec:nonce}), and tests of the paper's replay-resistance and forward-security claims
      (\cref{sec:claims}). \prov{Ours}, applying known techniques where stated.
\end{enumerate}

\paragraph{Organisation.}
\Cref{sec:background} gives the background and related work, \cref{sec:schemes} the two schemes, and
\cref{sec:method} our implementation and methodology. \Cref{sec:repro,sec:cost} cover the reproduction and
the cost audit; \cref{sec:keyonly,sec:opt,sec:agg,sec:nonce,sec:claims} the new investigations.
\Cref{sec:discussion} discusses what the results do and do not show, and \cref{sec:conclusion}
concludes. The code, tests, recorded measurements and the commands that reproduce every result are in
the project repository listed on the title page.
